\chapter{Team Experience and Market Evaluation}
\label{chapter:team-market}

%% ---------------------------------------------------------------
\section{Qualifications and Expertise}
\label{sec:qualifications}
%% ---------------------------------------------------------------

Emil Wes Lambert is completing a Bachelor of Science in Aerospace Engineering at TU Delft, Faculty of Aerospace Engineering. Relevant technical qualifications include:

\begin{itemize}
  \item \textbf{Spacecraft systems engineering}: SWAP budget analysis, hardware--software co-design, and system-level trade studies --- the methodology applied throughout this submission.
  \item \textbf{Attitude and orbit control systems (AOCS)}: The Kalman filter and MDP share the same Bayesian state estimation mathematical framework, making the MASFE scheduling policy a direct extension of AOCS-taught estimation theory.
  \item \textbf{Mission analysis and design}: Orbit mechanics, ground coverage analysis, and revisit scheduling --- directly applied in the orbital state representation of the MDP.
  \item \textbf{Embedded flight software}: Real-time operating systems, hardware--software partitioning, and resource-constrained computing --- the exact challenge of deploying the MDP policy on a LEON4FT + FPGA stack.
\end{itemize}

\noindent The agronomy and domain-expertise gap inherent in a solo aerospace submission is bridged through cited primary sources: Dr.\ Gold's Cornell spectral detection methodology~\cite{gold-webinar1}, USDA Economic Research Service crop loss data, and Loft Orbital's operational deployment experience~\cite{mytkowicz-webinar2}. MASFE is architected for single-operator mission management, validated by its solo development --- this lean operational concept reduces mission operations cost and reflects a commercial deployment philosophy where the orchestration engine is autonomous.

%% ---------------------------------------------------------------
\section{Past Projects and Relevant Research}
\label{sec:past-projects}
%% ---------------------------------------------------------------

Emil Wes Lambert's relevant project experience includes:

\begin{itemize}
  \item \textbf{Venus Atmosphere Mission Design (TU Delft DSE)}: Completed the Design Synthesis Exercise --- a 10-week full spacecraft design project delivering a balloon-based aerobot mission for long-duration Venus atmospheric and seismic investigation. Designed for 5$+$ year operational lifetime using in-situ nitrogen replenishment, within a \euro{}400\,M budget and 700\,kg mass constraint. Applied SWAP budgeting, subsystem trade-offs, and requirements tracing --- the same methodology used in this submission.
  \item \textbf{AeroDelft, Partnerships \& Business Manager (2 years)}: Served as Partnerships and Business Manager at AeroDelft, TU Delft's student team developing a hydrogen-powered aircraft. Led sponsorship negotiations, industry partnership development, and commercial market positioning --- directly applicable to MASFE's three-stream revenue model and MVP roadmap.
\end{itemize}

%% ---------------------------------------------------------------
\section{Scaling and Future Market Integration}
\label{sec:market}
%% ---------------------------------------------------------------

%% -----------------------------------
\subsection{Target Users and Design Criteria Validation}
\label{subsec:users}
%% -----------------------------------

Target users were identified and design criteria were validated through three channels:

\begin{enumerate}
  \item \textbf{Dr.\ Gold's Cornell research}~\cite{gold-webinar1} explicitly addresses the regenerative farming community's need for pre-symptomatic pathogen detection, providing both the spectral detection methodology and the user requirement for early-intervention alert timing.
  \item \textbf{OpenET's published FARMS API case studies}~\cite{openet-api} demonstrate active farmer demand for satellite-derived irrigation advisory products, validating the REST API delivery mechanism.
  \item \textbf{USDA Economic Research Service data} on precision agriculture adoption rates confirms that farm management platform integration (rather than a standalone app) is the correct delivery channel --- farmers already use platforms such as Climate FieldView and John Deere Operations Centre.
\end{enumerate}

\noindent Design criteria derived directly from user needs: the $\leq$90\,min latency requirement reflects the 24--48\,hour intervention window for fungicide application; the $\leq$30\,m GSD requirement reflects field-patch-level management practice; the REST API delivery requirement reflects the existing software ecosystem already deployed on commercial farms.

%% -----------------------------------
\subsection{Market Sizing}
\label{subsec:market-sizing}
%% -----------------------------------

\noindent\textbf{TAM/SAM/SOM.}
The global precision agriculture market is projected at $\sim$\$12\,B by 2028 (TAM). The satellite-derived crop analytics segment is $\sim$\$2.5\,B (SAM). Sub-day adaptive pathogen alerting represents a $<$\$50\,M greenfield segment today (SOM) --- no current product combines adaptive onboard MDP scheduling with multi-algorithm fusion for crop disease detection at field-patch scale. MASFE enters this uncontested position. Break-even at $\approx$500{,}000\,ha coverage is achievable with a three-satellite constellation at \$3--5/ha/year subscription pricing.

%% -----------------------------------
\subsection{Market Applications Beyond Agriculture and Forestry}
\label{subsec:beyond-ag}
%% -----------------------------------

MASFE's MDP orchestration engine is domain-agnostic. The scheduler, state representation, and reward function require only algorithm-specific parameter updates for new applications. Table~\ref{tab:generalisation} lists the primary generalisation pathways.

\begin{table}[htb]
  \centering
  \caption{MASFE generalisation beyond agriculture: the orchestration engine as a platform.}
  \label{tab:generalisation}
  \small
  \begin{tabularx}{\textwidth}{p{3.0cm}p{3.5cm}X}
    \toprule
    \textbf{Application} & \textbf{Swap algorithms to} & \textbf{Value proposition} \\
    \midrule
    Wildfire fuel moisture monitoring
      & FIRMS thermal anomaly + NDWI
      & Detect drying vegetation before ignition; adaptive revisit over high-risk terrain. Market: $\sim$\$600\,M by 2028, 12\,\% CAGR. \\
    \addlinespace
    Flood inundation mapping
      & SAR coherence change detection
      & MDP allocates revisit cadence over flood-front progression; inter-satellite cuing \\
    \addlinespace
    Urban heat island analytics
      & LST + albedo change
      & Same priority downlink logic for anomalous urban patches; city-scale coverage \\
    \addlinespace
    Defence / ISR
      & Anomalous activity detection
      & Same adaptive scheduling framework; replace stress index with activity score \\
    \bottomrule
  \end{tabularx}
\end{table}

\noindent The orchestration engine is the licensable IP. Agriculture is the first and primary market; the platform architecture serves any multi-algorithm Earth observation mission requiring intelligent onboard resource allocation.
